% !TEX root = ../bb_AiRa_Kappaleet.tex

% ö = \%c3\%b6
% ä = \%C3\%A4

\xdef\sectiontitle{\Isectiontitle} %aiheuttama
\TOCrefs

%%%%%%%%%%%%%%%
\begin{frame}[label=1_5_a]%[label=1_5_TOC1]
\frametitle{\texorpdfstring{\culine[\sectiontitleulinecolor]{\sectiontitle}}{\sectiontitle} }
\label{\TOCname}
\vspace{-0.5cm}

\begin{itemize}[leftmargin=0.5cm,itemsep=2mm]
    \item[1.1]<1-> \hyperlink{1_1_a}{\IaSubsectiontitle}
    \item[1.2]<1-> \hyperlink{1_2_a}{\IbSubsectiontitle}	
    \item[1.3]<1-> \hyperlink{1_3_a}{\IcSubsectiontitle}	
    \item[1.4]<1-> \hyperlink{1_4_a}{\IdSubsectiontitle}
    \item[1.5]<1-> \hyperlink{1_5_a}{\CDAlert[blue]{\IeSubsectiontitle}}
    \item[1.6]<1-> \hyperlink{1_6_a}{\IfSubsectiontitle}
    \item[1.7]<1-> \hyperlink{1_7_a}{\IgSubsectiontitle}
\end{itemize}

\vspace{-1cm}
\begin{itemize}[leftmargin=8.5cm,itemsep=2mm]
    \item[1.8]<1-> \hyperlink{1_8_a}{\IhSubsectiontitle}
	\item[1.9]<1-> \hyperlink{1_9_a}{\IiSubsectiontitle}
	\item[1.10]<1-> \hyperlink{1_10_a}{\IjSubsectiontitle}
	\item[1.11]<1-> \hyperlink{1_11_a}{\IkSubsectiontitle}
\end{itemize}

%\endofslide 
\end{frame}
%%%%%%%%%%%%%%%

%\xdef\IbSubsectiontitle{Entropy and Temperature}
\xdef\sectiontitle{\IeSubsectiontitle}

%% MAXWELL-BOLTZMANN
%%%%%%%%%%%%%%%
\begin{frame}%[label=1_5_a]
\frametitle{\texorpdfstring{\culine[\sectiontitleulinecolor]{\sectiontitle}}{\sectiontitle} }
\framesubtitle{Molekyylien vauhtijakauma}

\appear{}
\begin{itemize}
	\item Tarkastellaan seuraavaksi \CDAlert[red]{ideaalikaasusta koostuvaa systeemiä}\appear{, jossa yksittäisen molekyylin energia on}
\[
E  \equal \frac{1}{2} \appear{m} \appear{\textrm{((} v_{x}^{2}}\appear{+v_{y}^{2}}\appear{+v_{z}^{2} \textrm{))}}  \equal \frac{1}{2} \appear{m \appear{v^{2}}}
\]

\item \CDAlert[tunired]{Oletetaan että molekyylien nopeuksille on olemassa jakauma}\appear{, mutta emme vielä tiedä jakauman tarkkaa muotoa}

\item Löytääksemme mielivaltaisen \appear{keskinopeutta kuvaavan funktion $f=f(v)$}\appear{, meidän täytyy integroida 3D-tilavuudessa olevien molekyylien nopeuksien yli}\appear{, eli nopeusvektorilla on 3 komponenttia}

%\item The spatial integrations over volume $V$ cancel out, so we obtain:
%\[
%\overline{f(v)}=\frac{ \nint f(v) e^{-\frac{1}{2} m v^{2} / k_{B} T} d V d v_{x} d v_{y} d v_{z}}{ \nint e^{-\frac{1}{2} m v^{2} / k_{B} T} d V d v_{x} d v_{y} d v_{z}}=\frac{ \nint f(v) e^{-\frac{1}{2} m v^{2} / k_{B} T} d v_{x} d v_{y} d v_{z}}{ \nint e^{-\frac{1}{2} m v^{2} / k_{B} T} d v_{x} d v_{y} d v_{z}}
%\[
%Because we are searching for the speed distribution (not velocity), it is convenient to shift from cartesian to spherical coordinates using : $d v_{x} d v_{y} d v_{z}=v^{2} d v \sin \theta d \theta d  \phi $
\end{itemize}

\endofslide 
%\endofsection
\end{frame}
%%%%%%%%%%%%%%%

%%%%%%%%%%%%%%%
\begin{frame}
\frametitle{\texorpdfstring{\culine[\sectiontitleulinecolor]{\sectiontitle}}{\sectiontitle} }
\framesubtitle{Molekyylien vauhtijakauma}

\appear{}

\begin{itemize}
\item Integrointi tilavuuden $V$ yli sieventyy pois\appear{, jolloin saadaan:}
\[
\overline{f(v)}  \equal \fracc{ \appear{\nint} \appear{f(v)} \appear{e^{-\Frac{1}{2} \appear{m v^{2}} \appear{/ k_{B} T}}} \appear{d V'} \appear{d v_{x}} \appear{d v_{y}} \appear{d v_{z}}}{ \appear{\nint} \appear{e^{\appear{-\Frac{1}{2} m v^{2}} \appear{/ k_{B} T}}} \appear{d V' d v_{x} d v_{y} d v_{z}} }  \equal \fracc{ \nint \appear{f(v) e^{-\Frac{1}{2} m v^{2} / k_{B} T}} \appear{d v_{x} d v_{y} d v_{z}}}{ \nint e^{-\Frac{1}{2} m v^{2} / k_{B} T} d v_{x} d v_{y} d v_{z}}
\]
\item Koska haemme siis \CDAlert[blue]{vauhti}jakaumaa (ei \CDAlert[tuniblue]{nopeus})\appear{, on hyödyllistä \CDAlert[red]{siirtyä karteesisesta kordinaatistosta pallo} \CDAlert[red]{koordinaatistoon} käyttäen muunnosta: }
\[
d v_{x} d v_{y} d v_{z}  \equal \appear{v^{2} d v} \appear{\sin \theta \,d \theta} \appear{\,d \phi}
\]
\end{itemize}

\endofslide 
%\endofsection
\end{frame}
%%%%%%%%%%%%%%%


%%%%%%%%%%%%%%%
\begin{frame}
\frametitle{\texorpdfstring{\culine[\sectiontitleulinecolor]{\sectiontitle}}{\sectiontitle} }
\framesubtitle{Todennäköisyysjakauma}

\appear{}

\begin{itemize}
	\item Näin integraali muuttuu muotoon
	\[
	\overline{f(v)}  \equal  \frac{ \nint f(v) e^{-\Frac{1}{2} m v^{2} / k_{B} T} v^{2} d v \sin \theta d \theta d  \phi }{ \nint e^{-\Frac{1}{2} m v^{2} / k_{B} T} v^{2} d v \sin \theta d \theta d  \phi }
	 \equal  \frac{ \nint_{0}^{\infty} f(v) e^{-\Frac{1}{2} m v^{2} / k_{B} T} v^{2} d v}{ \nint_{0}^{\infty} e^{-\Frac{1}{2} m v^{2} / k_{B} T} v^{2} d v}
	\]
	% ALAKERTA ON: $=\sqrt{\frac{\pi}{2}}\left(\frac{k_{B} T}{m}\right)^{3 / 2}$
 \appear{missä nimittäjän integraalit}\appear{, ja muuttujien $\theta$} \appear{ja $\phi$}  \appear{yli tapahtuvat integraalit} \appear{voidaan ratkaista johtaen relaatioon}
\[
\overline{f(v)}  \equal  \appear{\nint_{0}^{\infty}  }\appear{f(v)} \appear{\LB} \appear{\sqrt{2/\pi}} \appear{\bigg( \Frac{m}{k_{B} T}  \bigg)^{3 / 2}} \appear{v^{2} e^{-\Frac{1}{2} m v^{2} / k_{B} T}}\appear{\RB} \appear{d v}
\]

\item Kaavan integrandin nähdään nyt muistuttavan \CDAlert[red]{todennäköisyysjakaumaa} $P=P(v)$:
\[
\overline{f(v)}  \equal  \appear{\nint f(v) d P} \equal  \appear{\nint f(v)} \frac{d P}{d v} \appear{d v}
\]

\end{itemize}

\endofslide 
%\endofsection
\end{frame}
%%%%%%%%%%%%%%%

%%%%%%%%%%%%%%%
\begin{frame}
\frametitle{\texorpdfstring{\culine[\sectiontitleulinecolor]{\sectiontitle}}{\sectiontitle} }
\framesubtitle{Maxwellin–Boltzmannin vauhtijakauma}

\appear{}

\begin{itemize}
\item Näin sulkujen sisällä oleva termi voidaan tulkita kaasumolekyylien vauhtien todennäköisyysjakaumaksi
\[
P(v)_{\text {Maxwell}}  \equal  \frac{d P}{d v}  \equal \appear{ \Sqrt{\Frac{2}{\pi}}   \bigg( \Frac{m}{k_{B} T} \bigg)^{3 / 2} v^{2} e^{-\Frac{1}{2} m v^{2} / k_{B} T}}
\]

\item Jos me \Alert[blue]{lisäämme molekyylien tilatiheyden $n=N / V$ kaavaan}\appear{, jakauman yksiköksi saadaan hiukkasta/yksikkötilavuus}\appear{, mutta jakauma sisältää informaatiota molekyylien vauhtien todennäköisyysjakaumasta}
\[
n(v)_{\text {Maxwell }}  \equal \frac{N}{V} \appear{\Sqrt{\Frac{2}{\pi}} } \appear{ \bigg( \Frac{m}{k_{B} T} \bigg)^{3 / 2}} \appear{v^{2}} \appear{e^{-\Frac{1}{2} m v^{2} / k_{B} T}}
\]
\appear{Tätä jakaumaa kutsutaan \Alert[tunired]{Maxwellin vauhtijakaumaksi}}\appear{, tai\\
\CDAlert[red]{Maxwellin–Boltzmannin vauhtijakaumaksi}}

\end{itemize}

\endofslide 
%\endofsection
\end{frame}
%%%%%%%%%%%%%%%


%%%%%%%%%%%%%%%
\begin{frame}
\frametitle{\texorpdfstring{\culine[\sectiontitleulinecolor]{\sectiontitle}}{\sectiontitle} }
\framesubtitle{Maxwellin–Boltzmannin vauhtijakauma}

\appear{\xdef\loc{north east}%
\begin{tikzpicture}[remember picture,overlay]% anchor=south
    \node (A) [yshift=\figYshift,xshift=\figXshift, anchor=\loc] at (current page.\loc) {\includegraphics[page=1,width=6cm]{Figures_booklet/SOM_9_Statistical_mechanics_figures/SOM_Figure_8.png}}; %scale=\figscale. % 8cm max hei
\end{tikzpicture}}%
% 6.5 cm = midway, 10 cm = 3/4 
\begin{minipage}{7.5cm}
	\begin{itemize}
\item Huomaa että vauhtijakauma kasvaa ylöspäin nollavauhdista (kuva~8)
	
%	\item Tämä ei tarkoita että \CDAlert[red]{todennäköisyys} ei kasvaisi energian kasvaessa! \appear{Todennäköisyys kasvaa}
\item Alkuun \CDAlert[red]{todennäköisyys} kasvaa tilan energian kasvaessa\appear{, koska jakauman neliöllinen nopeusriippuvuus hallitsee käytöstä}
%\item Tilojen lukumäärä \appear{(nopeusvektorit)} \appear{tie\-tyllä vauhdilla riippuu neliöllisesti vauhdista}

\item Korkeammilla tiloilla\appear{/nopeuksilla} \appear{mie\-hi\-tyksen toden\-näköi\-syys alkaa selkeästi}\appear{ \CDAlert[red]{vaimenemaan eksponentiaalisesti} \CDAlert[red]{energian funktiona}}

\end{itemize}
\end{minipage}
\vspace{3mm}

\begin{itemize}
\item[]<1->
$$
n(v)_{\text {Maxwell }}  = \Frac{N}{V} \Sqrt{\Frac{2}{\pi}}  \bigg( \Frac{m}{k_{B} T} \bigg)^{3 / 2} v^{2} e^{-\Frac{1}{2} m v^{2} / k_{B} T}
$$
\end{itemize}

%\endofslide 
\endofsection
\end{frame}
%%%%%%%%%%%%%%%


%%%%%%%%%%%%%%%%
%\begin{frame}
%\frametitle{\texorpdfstring{\culine[\sectiontitleulinecolor]{\sectiontitle}}{\sectiontitle} }
%\framesubtitle{test frame}
%
%\begin{itemize}
%\item testing frame:
%$$
%n(v)_{\text {Maxwell }}  = \Frac{N}{V} \Sqrt{\Frac{2}{\pi}}  \bigg( \Frac{m}{k_{B} T} \bigg)^{3 / 2} v^{2} e^{-\Frac{1}{2} m v^{2} / k_{B} T}
%$$
%and with left and right
%$$
%n(v)_{\text {Maxwell }}  = \Frac{N}{V} \Sqrt{\Frac{2}{\pi}}  \left( \Frac{m}{k_{B} T} \right)^{3 / 2} v^{2} e^{-\Frac{1}{2} m v^{2} / k_{B} T}
%$$
%\end{itemize}
%
%%\endofslide 
%\endofsection
%\end{frame}
%%%%%%%%%%%%%%%%
